\title{Byte Latent Transformer: Patches Scale Better Than Tokens}

\begin{abstract}
We propose the Byte Latent Transformer ({\textsc BLT}), a novel LLM framework operating directly on bytes, which achieves parity with traditional tokenization-based LLMs at scale while offering marked gains in both inference efficiency and robustness. {\textsc BLT} processes input by encoding bytes into variable-length patches, using these adaptive segments as the essential computational units. The boundary of each patch is determined by the entropy of the subsequent byte, thereby directing greater computational and model resources to regions with higher informational complexity. We conduct the first controlled study of byte-level model scaling at constant FLOPs, extending to architectures with up to 8B parameters and trained on 4T bytes. Our findings establish that large-scale learning from raw bytes—eschewing fixed vocabularies—is practical and scalable. Dynamic patch sizing enables better efficiency in training and inference, especially when longer patches are applied to more predictable data, and yields qualitative benefits in reasoning and coverage of rare cases. Taken together, for any fixed inference budget, {\textsc BLT} outperforms conventional tokenized models in scaling, benefitting jointly from the expansion of both patch lengths and model capacity.
\end{abstract}

\section{Introduction}
\label{section:intro}

We present the Byte Latent Transformer~(\textbf{BLT}), an architecture that operates directly on raw bytes without relying on tokenization, and for the first time demonstrates performance at par with token-based models when scaled, while achieving notable gains in both efficiency and robustness (\S\ref{sec:robustness}).
While most state-of-the-art large language models (LLMs) employ an almost entirely end-to-end learning framework, tokenization remains an exception: it is a heuristic pre-processing method that maps sequences of bytes into fixed sets of tokens. This process introduces inherent biases in how text sequences are represented, leading to issues such as sensitivity to domain or modality~\citep{dagangetting}, increased vulnerability to input perturbations~(\S\ref{sec:robustness}), inadequate encoding of orthographic information~\citep{edman2024cute}, and disparities across languages~\citep{liang2023xlm,petrov2024language,limisiewicz2024myte}.

Historically, tokenization has played a vital role because training llms directly on raw bytes is extremely resource-intensive at scale, mainly owing to increased sequence lengths~\citep{xue2022byt5}. To address these limitations, earlier studies have introduced more computationally efficient self-attention mechanisms~\citep{el2020characterbert, clark2022canine} or have opted for attention-free model designs~\citep{wang2024mambabyte}~(\S\ref{section:related_work}). Nevertheless, these improvements predominantly benefit the training of \textit{small models}. For large-scale models, most of the computational burden stems from the substantial feed-forward network layers that process each byte individually, rather than from the attention component itself.

For the purpose of optimizing computational resource distribution, we introduce a dynamic, learnable approach for aggregating bytes into \textit{patches}~(\S\ref{section:patching}), alongside a novel architecture that integrates both byte-level and patch-level information. In contrast to conventional tokenization, {\textsc BLT} dispenses with a fixed patch vocabulary; instead, it maps arbitrary byte groupings to latent patch representations through efficient, trainable encoder and decoder modules. Our findings indicate that this strategy enables a \textit{more} effective allocation of computational resources compared to models relying on tokenization.

Large language models that rely on tokenization process each token with an equal computational budget. This design prioritizes simplicity over optimal resource use, because tokenization often follows compression strategies that do not always align with the actual predictive difficulty of each token. Our model takes a different approach, emphasizing the flexible distribution of computation in response to task demands. For instance, predicting the final characters in a word typically requires less modeling capacity than generating the initial word in a new sentence, as the former task involves simpler, lower-entropy decisions. This principle shapes the design of {\textsc BLT}~(\S\ref{section:architecture}), which utilizes a hierarchical system of three transformer modules: two compact local byte-level models and one substantial global \textit{latent transformer}~(\autoref{fig:arch_high_level}). 

To allocate computation adaptively, {\textsc BLT} organizes bytes into patches according to the entropy involved in next-byte prediction, effectively grouping bytes together based on the similarity of their information content within context.

We conduct the first systematic scaling analysis of byte-level architectures regulated by floating point operations (flops), extending up to 8 billion parameters and processing 4 trillion bytes during training. Our results demonstrate the feasibility of fully end-to-end training directly on bytes, removing reliance on fixed-vocabulary tokenization. Notably, {\textsc BLT} achieves comparable performance to Llama 3 in terms of training efficiency when normalized by flop consumption,\footnote{We determine a model’s computational expenditure by tallying the number of floating point operations (flops) executed.} while reducing inference flops by as much as 50\%~(\S\ref{section:scaling}).

Furthermore, we find that modeling directly with byte sequences confers substantial gains in handling rare or long-tail distributions. {\textsc BLT} models exhibit superior robustness to noisy input and improved capabilities in character-level phenomena, as demonstrated across orthography, phonological patterning, and low-resource machine translation tasks~(\S\ref{sec:robustness}).

Additionally, our investigations reveal that {\textsc BLT} allows us to scale both model and patch sizes concurrently, while keeping the inference flop budget fixed. Increasing patch lengths typically reduces computation requirements, enabling the allocation of computational resources towards expanding the global latent transformer component, which executes less frequently. Our inference-flop-normalized scaling studies~(\autoref{fig:fixed_inference_scaling}) highlight marked improvements in scaling characteristics compared to models that depend on token-based representations.

To conclude, the main contributions of this paper are as follows: 1) The proposal of {\textsc BLT}, a byte latent llm framework that adaptively distributes computation for enhanced flop efficiency; 2) Empirical evidence that our approach matches the training flop requirements of Llama 3 models up to the 8B parameter scale, with the added flexibility to achieve up to 50\% flop efficiency improvements at the expense of only slight decreases in evaluation performance; 3) The capability of {\textsc BLT} models to facilitate new scalability strategies for llms, enabling growth in model size without surpassing predetermined inference cost constraints; 4) Demonstration that {\textsc BLT} models exhibit increased robustness against noisy inputs and greater sensitivity to sub-word characteristics that typically elude token-level llms. All training and inference resources for {\textsc BLT} are made available at~\url{https://github.com/facebookresearch/blt}.

\section{Patching: From Individual Bytes to Groups of Bytes}
\label{section:patching}

Dividing bytes into \textit{patches} enables {\textsc BLT} to assign computational resources adaptively, taking contextual information into account. As illustrated in Figure~\ref{fig:patching_types}, there are multiple approaches to grouping bytes into patches. More precisely, a patching function $f_p$ takes a byte sequence $\pmb{x}=\{x_i,|i=1,\ldots n\}$ of length $n$ and divides it into a shorter sequence of $m < n$ patches $\pmb{p}=\{p_j|j=1,\ldots,m\}$, where each $x_i$ is mapped to \{0,1\}; the value 1 signals the initiation of a new patch. In both token-oriented and patch-oriented frameworks, the computational requirements largely depend on how many steps the primary Transformer must perform. For {\textsc BLT}, this equates to the total patches formed by a specific patching function. Thus, the mean patch length, often referred to as \textit{patch size}, is crucial in assessing the compute demands for both training and inference phases under a chosen patching function~(\S\ref{section:flops}). 

We now present three kinds of patching functions: segmentation based on a constant byte count per patch~(\S\ref{section:static-patch}), segmentation at whitespace boundaries~(\S\ref{section:white-patch}), and adaptive patching based on entropy values produced by a compact byte language model~(\S\ref{section:dyn-patch}). We further elaborate on the concept of incremental patching and contrast patching with conventional tokenization~(\S\ref{section:bpe-patch}).

\subsection{Strided Patching Every K Bytes}
\label{section:static-patch}

One of the simplest methods for organizing bytes is to aggregate them into fixed-length patches of size $k$, as implemented in MegaByte~\citep{yu2023megabyte}. This uniform stride is straightforward to apply during both training and inference, allows for a clear approach to adjusting the average patch length, and thus facilitates direct control over computational costs. Nevertheless, there are notable limitations to this approach. Primarily, the allocation of computational resources does not adapt to content: transformer steps can be wasted when dealing with less informative segments (such as sequences of whitespace in code), or alternatively, patches with complex content—like mathematical expressions—may not receive adequate computation. Additionally, this static partitioning can cause identical or similar byte sequences to be segmented differently, resulting in inconsistent and context-insensitive patching—for example, the same word might be split in various ways across occurrences.

\subsection{Space Patching}
\label{section:white-patch}

According to \citet{slagle2024spacebyte}, an effective enhancement to strided patching is introduced, whereby new patches are generated after the appearance of any space-like bytes\footnote{
    Space-like bytes are categorized as any byte that does not represent a latin letter, numeric digit, or utf-8 continuation byte. Furthermore, it is required that each patch includes at least one byte that is not space-like. }. This approach capitalizes on the fact that space-like bytes often delineate the boundaries of linguistic units in various languages. The Space patching method allots a transformer latent computation step (i.e., increased computational effort) to model each word individually. As a result, word-level patching remains consistent throughout sequences, and greater computational resources are dedicated to segments likely to demand challenging predictions—commonly those that occur after spaces. For instance, the prediction difficulty of the initial byte of a response to “Who composed the Magic Flute?\uline{\hspace{.6em}}” exceeds that of subsequent bytes following the character “M”. This is because the initial letter drastically narrows down the plausible options, simplifying the remainder of the answer, such as in the completion “Mozart”. Nonetheless, Space patching exhibits limitations, particularly in its inability to adapt gracefully across different languages and application areas, and more crucially, its fixed patch size. In the subsequent section, we describe an alternative patching strategy inspired by the observation that initial bytes in words usually pose the greatest prediction challenge, while also affording a flexible approach to patch size selection.

\subsection{Entropy Patching: Using Next-Byte Entropies from a Small Byte LM}
\label{section:dyn-patch}

Instead of utilizing a rule-based strategy like whitespace, we adopt a data-driven method to pinpoint points of elevated next-byte prediction uncertainty. Specifically, we propose \textit{entropy patching}, a technique that leverages entropy measurements to determine the locations of patch boundaries.

We train a small byte-level auto-regressive language model on the training data for {\textsc BLT} and compute next byte entropies under the LM distribution $p_{e}$ over the byte vocabulary $\mathcal{V}$:

\begin{equation}
H(x_i) = \sum_{v \in \mathcal{V}} p_{e}(x_i=v|\pmb{x}_{<i}) \log p_{e}(x_i=v|\pmb{x}_{<i})
\end{equation}

We investigate two approaches for detecting patch boundaries using the entropy values $H(x_i)$. The first method selects points whose entropies surpass a predetermined global threshold, as shown in~\autoref{fig:patching}. The second method highlights locations where the entropy is notably higher compared to the preceding value. This latter technique may also be viewed as pinpointing positions where there is a disruption in the generally monotonically decreasing entropy within a patch.

\begin{align*}
    \text{Global Constraint}\hspace{1cm}H(x_t)& > \theta_g\\
    \text{Approx. Monotonic Constraint}\hspace{1cm}H(x_t)& - H(x_{t-1}) > \theta_r
\end{align*}

Patch boundaries are determined through a simple preprocessing procedure that occurs as part of the dataloading process. This approach contrasts with the method used by~\citet{nawrot-etal-2023-efficient}, where a classifier is employed to learn and predict patch boundaries based on entropy measures. In our experimental analysis~(\S\ref{section:experiments}), we conduct a comparison between these two techniques for separating bytes with low entropy from those with high entropy.

\subsection{The Byte-Pair Encoding (BPE) Tokenizer and Incremental Patching}
\label{section:incremental-patching}
\label{section:bpe-patch}

A variety of contemporary llms, such as our reference model Llama 3, employ subword tokenization methods like bpe~\citep{gage1994new, sennrich-etal-2016-neural}. In our terminology, ``tokens'' denote byte sequences selected from a \textit{finite} set established before training begins, whereas ``patches'' indicate adaptively grouped segments that do not rely on a pre-defined vocabulary. One fundamental distinction between these is that tokens prevent the model from directly accessing the original byte-level information.

An important advancement offered by {\textsc BLT} relative to tokenization-based approaches is its reconceptualization of the balance between vocabulary magnitude and computational requirements. Traditional language models operate under the constraint that enlarging the vocabulary leads to an increase in the average token length, thereby reducing the total number of processing steps; however, it also expands the dimensionality of the model's final output projection layer. Consequently, this relationship severely restricts how much tokenization-based methods can modulate both token length and inference expenses. For instance, Llama 3 achieves a rise in average token length from 3.7 to 4.4 bytes, but this improvement comes with a fourfold enlargement of its embedding matrix relative to Llama 2.

During generation, {\textsc BLT} must determine at each step of the byte sequence whether it is positioned at a patch boundary. This identification governs if the Latent Transformer should be engaged for additional computation. Notably, this determination must be made without reference to future bytes that have not yet been generated; in other words, patch segmentation decisions must be made solely based on the existing prefix. More formally, we say that a patching scheme $f_p$ supports incremental patching if the following condition holds:
$$f_p(\pmb{x}_{<i}) = f_p(\pmb{x})_{<i}$$
bpe, however, fails to constitute an incremental patching scheme, as the way a prefix is tokenized can vary depending on what follows in the sequence, thus it does not fulfill the above property\footnote{Introducing a designated delimiter token to signal patch boundaries would make bpe incremental in this regard, albeit at the cost of longer byte sequences.}.

\section{{\textsc BLT} Architecture}
\label{section:architecture}
{\textsc BLT} consists of a large-scale, global autoregressive language model designed to process patch representations. In addition, two more compact local models are used: one to transform byte sequences into patches, and another to reconstruct bytes from patch representations~(\autoref{fig:arch_high_level}).

\subsection{Latent Global Transformer Model}

\textit{The Latent Global Transformer} refers to an autoregressive transformer model, denoted $\mathcal{G}$, comprised of $l_{\mathcal{G}}$ layers. This model transforms a series of input latent patch representations, $p_j$, into corresponding output patch representations, $o_j$. In this manuscript, we employ the subscript $j$ for referencing patches and $i$ for bytes. The global model incorporates a block-causal attention mechanism~\citep{dubey2024llama}, which confines attention to the present and preceding patches within the current document. Notably, this model accounts for the majority of computational operations (flops) during both pre-training and inference stages. By selectively determining when the global transformer is applied, one can regulate and modulate computational resources allocated to distinct input and output segments, in accordance with the respective complexity of those segments.

\subsection{Local Encoder}
\label{section:local}

The \textit{Local Encoder Model}, represented by $\mathcal{E}$, employs a transformer architecture designed to be lightweight, comprising significantly fewer layers ($l_{\mathcal{E}} << l_{\mathcal{G}}$) compared to the global transformer. Its central purpose is to rapidly translate sequences of input bytes, $b_i$, into rich patch-level representations, $p_j$. In contrast to standard transformer models, this encoder incorporates cross-attention modules following every transformer layer, enabling aggregation of byte-level features into patch-level representations as illustrated in~(\autoref{fig:crossattn}).

Initially, input bytes $b_i$ are mapped into a vector space using an embedding table of shape $\mathbb{R}^{256 \times h_{\mathcal{E}}}$, yielding $x_i$. Optionally, these representations may be supplemented with hash-embeddings~(\S\ref{sec:hash-emb}) to encode extra information.

Sequences of alternating transformer and cross-attention modules~(\S\ref{sec:enc-cross-attn}) act on these embeddings, culminating in patch representations, $p_i$, which serve as inputs for the global transformer $\mathcal{G}$.

Within this model, transformer layers make use of a \textit{local block causal} attention masking strategy: each byte’s attention window is constrained to a fixed number $w_{\mathcal{E}}$ of prior bytes. This windowing permits attention to cross patch boundaries when necessary but enforces strict containment within document boundaries. Subsequent subsections elaborate on specifics pertaining to the embedding scheme and the architecture of the cross-attention block.

\subsubsection{Encoder Hash n-gram Embeddings}
\label{sec:ngram-emb}
\label{sec:hash-emb}
To form effective and nuanced representations at each position $i$, it is essential to account for contextual information from prior bytes. In {\textsc BLT}, this is addressed by encoding the byte $b_i$ not only in isolation, but also within the context of a surrounding byte n-gram. Specifically, at each timestep $i$, we begin by generating byte-grams

\begin{align}
    g_{i,n}=\{b_{i-n + 1},\ldots, b_i\}
\end{align}

for each byte position $i$ and for values of $n$ ranging from three up to eight.\footnote{
Byte-grams of length $n$ or greater are omitted whenever $i < n$. }

Next, we define hash $n$-gram embeddings, where every byte $n$-gram is mapped through a hash function to an entry in a corresponding embedding table $E_{n}^{hash}$ of predetermined size, specific to each $n \in\{3,4,5,6,7,8\}$~\citep{bai2010learning}.
This computed embedding is combined additively with the embedding for the byte itself, normalized, and subsequently fed as input to the local encoder. The resultant enhanced embedding is given by

\begin{align}
e_i &= x_i + \sum_{n = 3,...,8}E_{n}^{hash}(\text{Hash}(g_{i,n})) \\
\text{where, Hash}(g_{i,n}) &= \text{RollPolyHash}(g_{i,n}) \% |E_{n}^{hash}|
\end{align}

Each $e_i$ is scaled by dividing by the total number of considered $n$-gram sizes incremented by one, utilizing the RollPolyHash function as described in Appendix~\ref{appendix:rollpolyhash}. Section~\ref{section:ablations} investigates the influence of varying both the values of $n$ and the embedding table size for $n$-gram hash embeddings on flop-controlled scaling law behavior. Beyond hash-based $n$-gram embeddings, we also assessed embeddings based on $n$-gram frequencies; further information about these experiments is available in Appendix~\ref{appendix:ngrams}.

\subsubsection{Encoder Multi-Headed Cross-Attention}
\label{sec:enc-cross-attn}
Our approach is modeled after the cross-attention mechanism for inputs introduced in the Perceiver architecture~\citep{jaegle2021perceiver}. The key distinction in our method lies in the fact that, instead of employing a predetermined set of latent representations, we assign latent vectors to the variable patch representations~(\autoref{fig:crossattn}), and restrict attention to only those bytes that are part of the associated patch. Within the module, each patch $p_j$ is represented by its own query vector, which we obtain by aggregating (pooling) the corresponding byte-level representations for $p_j$. This pooled vector is then mapped through a linear transformation, $\mathcal{E}_{C} \in \mathbb{R}^{h_{\mathcal{E}} \times (h_{\mathcal{E}} \times U_{\mathcal{E}})}$, where $U_{\mathcal{E}}$ denotes the number of cross-attention heads in the encoder. Explicitly, letting $f_{\text{bytes}}(p_j)$ be the sequence of bytes belonging to patch $p_j$, we proceed to compute

\begin{align}
P_{0,j} &= \mathcal{E}_{C}(f_\text{bytes}((p_j)), f~ \text{is a pooling function} \\
P_l &= P_{l-1} + W_o\left(\text{softmax}\left(\frac{QK^T}{\sqrt{d_k}}\right)V\right) \\
\text{where } Q_j &= W_q(P_{l-1,j}), K_i = W_k(h_{l-1,i}), V_i = W_v(h_{l-1,i}) \\
h_l &= \text{Encoder-Transformer-Layer}_l(h_{l-1})
\end{align}

Here, $P \in \mathbb{R}^{n_p \times h_{\mathcal{G}}}$ denotes the representations for $n_p$ patches intended for input to the global model. These patch representations are produced by aggregating the byte embeddings $e_i$ associated with each patch $p_j$. The matrices $W_q$, $W_k$, $W_v$, and $W_o$ serve as the linear transformations that generate queries, keys, values, and outputs, respectively. In this setting, the keys and values are derived from byte-level representations $h_i$ produced by the preceding layer (or from $e_i$ at the first layer). The masking scheme employed for patching ensures that each query $Q_j$ attends exclusively to keys and values derived from bytes within its corresponding patch $j$. To accommodate the higher dimensionality commonly found in patch representations ($h_{\mathcal{G}}$ versus $h_{\mathcal{E}}$), the multi-headed attention is implemented such that $P_l$ is composed of multiple heads, each with dimension $h_{\mathcal{E}}$ during cross-attention; these multiple head outputs are subsequently concatenated to recover the original $h_{\mathcal{G}}$ dimensionality. A pre-LayerNorm is applied to the queries, keys, and values before the attention operation, and this cross-attention component does not employ any positional embeddings. The entire cross-attention mechanism is finally enclosed by a residual (skip) connection.

\subsection{Local Decoder} 

Analogous to the local encoder, the local decoder $\mathcal{D}$ utilizes a compact transformer architecture, featuring $l_{\mathcal{D}} << l_{\mathcal{G}}$ layers, and is designed to reconstruct raw bytes $y_i$ from a series of global patch embeddings $o_j$. This decoder generates each raw byte sequentially, conditioned on all previously produced bytes, leveraging the latent representations generated by the local encoder for the input byte string. Its architecture consists of $l_{\mathcal{D}}$ repeated blocks, each comprising a cross-attention module followed by a transformer layer. Here, the decoder’s cross-attention module precedes its transformer block to facilitate the conversion of patch-level features into byte-specific representations, after which the local decoder’s transformer module processes the resultant sequence of bytes.

\subsubsection{Decoder Multi-headed Cross-Attention} 

In the decoder cross-attention, the roles of the queries and key/values are interchanged i.e. the byte-representations are now the queries, and the patch representations are now the key/values. The initial byte-representations for the cross-attention are initialized as the byte embeddings from the last encoder layer i.e. $h_{l_{\mathcal{E}}}$. The subsequent byte-representations for layer $l$, $d_{l,i}$ are computed as:

\begin{align}
D_0 &= h_{l_{\mathcal{E}}} \\
B_l &= D_{l-1} + W_o\left(\text{softmax}\left(\frac{QK^T}{\sqrt{d_k}}\right)V\right), \\
  &\text{where } Q_i = W_q(d_{l-1,i}), K_i = W_k(\mathcal{D}_C(o_j)), V_i = W_v(\mathcal{D}_C(o_j)) \\
  D_l &= \text{Decoder-Transformer-layer}_l(B_l)
\end{align}

Here, $W_k$ and $W_v$ serve as projection matrices for keys and values, applying both a linear transformation and the split operation $\mathcal{D}_C$ to the final patch representations $o_j$ derived from the global model. Meanwhile, $W_q$ denotes the query projection matrix applied to byte representations $d_{l-1}$ from the preceding decoder transformer layer, or to $h_{l_{\mathcal{E}}}$ in the initial layer. $W_o$ refers to the output projection matrix, resulting in $B \in \mathbb{R}^{h_{\mathcal{D}} \times n_b}$, where $n_b$ is defined as the number of output bytes. The subsequent decoder states $D_l$ are obtained by passing the output of the cross-attention module, $B$, through a decoder transformer layer. Following the approach used in the local encoder cross-attention mechanism, we employ multiple attention heads, implement pre LayerNorms, omit positional embeddings, and utilize a residual connection encircling the cross-attention component.

\section{Experimental Setup}
\label{section:experiments}
We structured controlled experiments with the aim of rigorously comparing {\textsc BLT} and models relying on tokenization, ensuring that {\textsc BLT} receives no unfair benefit from the potential to exploit longer contextual sequences.

\subsection{Pre-training Datasets} In this study, we pre-train all model sizes using two primary datasets: (1) The Llama 2 dataset~\citep{touvron2023llama}, consisting of 2 trillion tokens sourced from numerous public datasets that have been subsequently curated and refined to enhance data quality; and (2) {\textsc BLT}-1T, a newly assembled dataset of 1 trillion tokens drawn from diverse public sources, which also incorporates selected portions of the Datacomp-LM pre-training release~\citep{li2024datacomp}. The Llama 2 dataset serves as the basis for scaling law analyses, following the optimal token quantity recommendations in~\cite{dubey2024llama} to guide the design decisions for {\textsc BLT}. In contrast, the {\textsc BLT}-1T dataset is utilized for comprehensive pre-training in order to enable a head-to-head comparison with Llama 3 across a range of downstream benchmarks. Notably, neither dataset incorporates content derived from Meta products or services. Additionally, in establishing baseline performances for tokenizer-based models, we opt for the Llama 3 tokenizer, which comprises a 128K token vocabulary; this choice outperformed the Llama 2 tokenizer in our preliminary experiments.

\subsection{Entropy Model} In our experiments, the entropy model employed is a byte-level language model that is trained on the same data distribution as the full {\textsc BLT} model. Except where noted otherwise, our default configuration utilizes a transformer architecture composed of 14 layers, a model size of 100M parameters, and a hidden dimension of 512, along with a sliding window attention mechanism that covers 512 bytes. All additional hyperparameter settings are aligned with those used in the local and global transformer models. We have also evaluated the impact of varying model capacity, receptive field size, and architectural choices, as outlined in \autoref{section:ablations}. Notably, for sufficiently restricted receptive fields, the trained entropy model can be represented compactly using a lookup table.

\subsection{Entropy Threshold and Equalizing Context Length} For approaches that incorporate entropy-based patching, we determine a suitable patching threshold intended to yield a targeted mean \textit{patch size} across the pretraining data blend. Within {\textsc BLT}, as opposed to tokenization, the selection of \textit{patch size} is flexible and can be set arbitrarily, which substantially influences the context length the model utilizes. To prevent models with larger patch sizes from obtaining an undue benefit via extended context, we hold the expected byte count per batch fixed. Accordingly, models configured with larger patch sizes are assigned shorter sequence lengths. Specifically, when using Llama 2 data, each context is set to 8k bytes, whereas for the {\textsc BLT}-1T dataset, the average context is increased to 16k bytes, while retaining an average batch size of 16M bytes.

Although the average batch size remains fixed, dynamic patching approaches produce varying proportions of bytes to patches in each batch. In our approach, the implementation of {\textsc BLT} training assembles batches by aggregating patches, thereby preventing the need for padding operations within the computationally intensive latent transformer. This method maintains a consistent number of patches per batch. Throughout training, we adjust byte sequences—padding or truncating as necessary—to 12k bytes for Llama 2 and 24k bytes for {\textsc BLT}-1T datasets. This strategy helps mitigate memory surges caused by sequences containing exceptionally large patches.

\subsection{Entropy Model Context}
\label{section:entropy-context}
Our empirical results indicate that entropy patching produces increasingly larger patches in contexts characterized by structured data, such as multiple choice tasks, which frequently exhibit a high degree of repetitiveness (refer to the MMLU patching demonstration in \autoref{fig:mmlu_patching}). The underlying cause of this phenomenon is the reduced entropy associated with the repeated information in the entropy model context. Consequently, for our large-scale experiment with {\textsc BLT}-Entropy at patch size 4.5, we refresh the entropy context by introducing new lines and implement an approximate monotonicity constraint, as this approach exhibits a lower susceptibility to "entropy drift" resulting from variations in context length. Notably, this adjustment is limited to the manner in which entropies are calculated; the method for determining the entropy threshold remains unchanged.

\subsection{FLOPs Estimation} 
\label{section:flops}

Our approach to calculating transformer FLOPs closely mirrors the methodology outlined in Chinchilla~\citep{hoffmann2022training}, which accounts for FLOPs originating from feed-forward modules, qkvo projections in the self-attention mechanism, as well as attention calculation and output projection. One important distinction is our assumption that the input embedding layer is realized through an efficient lookup table, rendering its computational cost negligible (effectively 0 FLOPs), rather than as a dense matrix multiplication. Consistent with established practices, we further estimate the FLOPs required for the backward pass to be double that of the forward pass.

To compute flops \textit{per byte} for {\textsc BLT} models, we add up the flops for the local encoder transformer, the global latent transformer, and the local decoder transformer, together with the cross attention blocks in the encoder and the decoder:

\begin{align}
    \text{FL}_{\text{{\textsc BLT}}} &= \text{Transf. FL}(h_{\mathcal{G}}, l_{\mathcal{G}}, m=n_{ctx}/n_p, V=0) / n_p\\
   &+\text{Transf. FL}(h_{\mathcal{E}}, l_{\mathcal{E}}, m=w_{\mathcal{E}}, V=0) \\
   &+ \text{Transf. FL}(h_{\mathcal{D}}, l_{\mathcal{D}}, m=w_{\mathcal{D}}, V=256) \\ 
    &+ \text{Cross Attn. FL}(h_{\mathcal{E}}, l_{\mathcal{E}}, m=n_p, r=n_p/k) \cdot k/n_p \\ 
   &+ \text{Cross Attn. FL}(h_{\mathcal{D}}, l_{\mathcal{D}}, m=k, r=k/n_p)
\end{align}

In this context, $n_{ctx}$ denotes the length of the sequence measured in bytes, while $n_p$ refers to the patch size. The parameter $r$ represents the proportion of queries relative to keys/values, and $k$ signifies the ratio of the patch dimension to the byte dimension—meaning the total count of local model segments that are concatenated to construct the overall model representation (as illustrated by $k=2$ in \autoref{fig:crossattn}). The variable $V$ stands for the output projection vocabulary size, which is exclusively utilized in the local decoder. The value of $m$—used to determine the context length for attention operations—depends on whether the respective module operates over bytes or patches. Accordingly, we adjust the calculation of attention-related flops for each module. Detailed formulations for the calculation of Transformer-FLOPs as well as Cross-Attention FLOPs are included in Appendix~\ref{section:flops-appendix}.

\subsection{Bits-Per-Byte Estimation} Perplexity only makes sense in the context of a fixed tokenizer as it is a measure of the uncertainty for each token. When comparing byte and token-level models, following previous work~\citep{xue2022byt5, yu2023megabyte, wang2024mambabyte}, we instead report Bits-Per-Byte (BPB), a tokenizer independent version of perplexity. Specifically:

\begin{align}
    \text{BPB}(x)&= \frac{\mathcal{L}_{CE}(\pmb{x})}{\ln(2)\cdot n_{\text{bytes}}}
\end{align}

Here, the cross-entropy loss accumulated over the data $\pmb{x}$ is divided by both the total byte count of $\pmb{x}$ and a fixed scaling constant to provide a normalized measure of uncertainty.

\subsection{Transformer Architecture Hyperparameters} The transformer blocks utilized in {\textsc BLT}, encompassing both the local and global models, are primarily based on the Llama~3~\citep{dubey2024llama} framework. Within the feed-forward modules, we adopt the SwiGLU activation~\citep{swiglu}, and for self-attention components, rotary positional embeddings (RoPE)~\citep{RoPe} parameterized by $\theta=500000$~\citep{xiong2024effective} are applied exclusively. RMSNorm \citep{RMSNorm} is employed for normalizing layers throughout. For the self-attention layers that employ standard, unchanging attention masks such as \textit{block causal} or \textit{fixed-window block causal}, we implement Flash attention~\citep{dao2022flashattention}, specifying a window size of 512 in the fixed-width scenario. In contrast, cross-attention modules incorporate dynamically-generated masks contingent on patch structure, for which we rely on Flex Attention\footnote{\url{https://pytorch.org/blog/flexattention}} to enable fused executions and substantially accelerate the training process.

\subsection{{\textsc BLT}-Specific Hyperparameters}

To evaluate the performance of {\textsc BLT} models, our experimental design explores two main axes: examining scaling properties and measuring results on various downstream tasks. For these analyses, we investigate models across a range of parameter sizes—specifically, 400M, 1B, 2B, 4B, and 8B. Architectural settings for these configurations can be found in Appendix~Table~\ref{tab:arch}.

The queries for the initial cross-attention layer of the local encoder are initialized using max-pooling. For every {\textsc BLT} model, we utilize $500,000$ hashes with one hash function, and n-gram sizes spanning from 3 to 8. We employ a fixed learning rate of $4e-4$ throughout. The decision to align the learning rates between the token baseline and {\textsc BLT} variants was based on hyperparameter searches at 400M and 1B scales, varying rates between $1e-3$ and $1e-4$, which indicated an identical optimal rate for both. When examining scaling effects on Llama-2 training data, batch sizes correspond to recommendations in~\cite{dubey2024llama}, or to matching byte-equivalent values. For optimization, AdamW~\citep{adamw} is adopted with $\beta_1 = 0.9$, $\beta_2 = 0.95$, and $\epsilon=10^{-8}$. The learning rate undergoes a linear warm-up for 2000 steps, followed by cosine decay down to zero. Additionally, a weight decay of 0.1 and a global gradient clipping threshold of 1.0 are enforced.

\section{Scaling Trends}
\label{section:scaling}

We provide a comprehensive overview of the scaling dynamics observed in byte-level models, offering insights that can support the continued development of {\textsc BLT} models. Our scaling investigation seeks to overcome shortcomings identified in earlier studies of byte-level modeling by: (a) systematically examining scaling behavior under compute-optimal training setups, (b) training parallel 8B-parameter models with substantial datasets (reaching up to 1T tokens or 4T bytes) and conducting evaluations on downstream benchmarks, and (c) analyzing scaling laws under scenarios where inference cost is kept constant. In a subsequent section, we will explore in detail the distinct benefits that arise from modeling sequences at the byte level.

\subsection{Parameter Matched Compute Optimal Scaling Trends}
\label{section:parameter-matched-scaling}
Leveraging the Llama 2 dataset, we train both \textit{compute-optimal} bpe models and {\textsc BLT} models at four distinct model capacities, spanning from 1B to 8B parameters. Training flops are then plotted in relation to language modeling accuracy on a representative sample of the training data composition. The bpe models employ the model parameter-to-dataset size ratio identified as optimal by Llama~3~\citep{dubey2024llama}, with this \textit{compute-optimal} arrangement theoretically delivering peak performance for a specified compute allocation~\citep{hoffmann2022training}, thus serving as a reliable reference point. Alongside each bpe model, a {\textsc BLT} counterpart is trained on identical data using a Latent Transformer precisely matched in architectural design and parameter count.

As shown in~\autoref{fig:scaling-compute-opt} (right), {\textsc BLT} models consistently achieve performance that is comparable to or exceeds that of bpe models, a pattern that persists as both model parameters and computational resources are increased. According to our knowledge, {\textsc BLT} represents the first byte-level Transformer to demonstrate comparable scaling behavior to BPE-based models under compute-optimal conditions. This finding substantiates our hypothesis that the parameter-to-compute ratio optimal for bpe models is also applicable to {\textsc BLT}, or at least not significantly different.

Both enhancements to architecture and the use of dynamic patching are essential for achieving parity with bpe scaling patterns. In ~\autoref{fig:scaling-compute-opt} (left), we evaluate models utilizing space-patching against Llama 3, estimating SpaceByte \citep{slagle2024spacebyte} by employing {\textsc BLT} space-patching while omitting n-gram embeddings and cross-attention mechanisms. While SpaceByte demonstrates gains relative to Megabyte, there remains a notable performance gap compared to Llama 3. As shown in \autoref{fig:scaling-compute-opt} (right), both structural adjustments and dynamic patching drive additional performance gains. At scale, {\textsc BLT} models achieve results comparable to leading tokenizer-based approaches such as Llama 3.

Additionally, we note that the selection of tokenizer has a noticeable impact on the performance of tokenizer-based models; specifically, models utilizing the Llama-3 tokenizer achieve better results compared to those employing the Llama-2 tokenizer when trained on identical datasets.

Lastly, our {\textsc BLT} model displays an intermediate behavior relative to Llama 2 and 3 when patch sizes are substantially increased. While the average token lengths for the Llama 2 and Llama 3 bpe tokenizers are 3.7 and 4.4 bytes, respectively, {\textsc BLT} demonstrates comparable scaling characteristics at average patch sizes of 6 or even 8 bytes. Since inference flop decreases as the average patch size increases, opting for an 8-byte patch size can yield nearly a 50\% reduction in inference flops. Additionally, models with larger patches show improved performance as both model scale and dataset size increase. Notably, {\textsc BLT} at an 8-byte patch size starts with inferior performance relative to bpe-based Llama 2 at 1B, but surpasses it at the 7B scale. This pattern indicates that using such large patch sizes could become even more effective as models and training budgets continue to expand, and that employing even greater patch sizes may be plausible with further scaling.

\subsection{Beyond Compute Optimal Task Evaluations}
\label{section:evals}
To investigate scaling behaviors more comprehensively, we train an 8B {\textsc BLT} model that exceeds the compute optimal regime, utilizing the {\textsc BLT}-1T dataset—a larger and higher-quality data source. We then evaluate model performance using a range of established benchmarks covering both classification and generation. For this assessment, we include tasks related to common sense reasoning, world knowledge, and code generation:
\paragraph{Classification tasks} comprise ARC-Easy (0-shot)~\citep{Eval_ARC}, Arc-Challenge (0-shot)~\citep{Eval_ARC}, HellaSwag (0-shot)~\citep{Eval_hellaswag}, PIQA (0-shot)~\citep{Eval_piqa}, and MMLU (5-shot)~\citep{hendrycks2020measuring}.  We utilize prompt-based scoring, computing the probabilities assigned to option tokens, and present the mean accuracy across tasks.
\paragraph{Coding related generation tasks:} We evaluate the model’s ability to write Python code using pass@1 metrics on MBPP (3-shot)~\citep{Eval_MBPP} and HumanEval (0-shot)~\citep{Eval_HumanEval}.

In \autoref{tab:evals}, we present a comparison among three models, each trained on the {\textsc BLT}-1T dataset: a baseline model utilizing the Llama 3 tokenizer and byte pair encoding (bpe),\footnote{The Llama 3 tokenizer, featuring a 128k vocabulary, was selected due to its superior performance compared to Llama 2's 32k vocabulary.} alongside two distinct configurations of the {\textsc BLT} model. The first variant adopts a space-patching strategy ({\textsc BLT}-Space), while the second employs a patching mechanism grounded in entropy ({\textsc BLT}-Entropy). For the entropy model, we enforce an approximate monotonicity constraint and reset its context at each new line, following the method described in~\autoref{section:entropy-context}. Training for all three models was performed with a matching flop allocation. Additionally, for {\textsc BLT}-Entropy, we decrease the entropy threshold during inference from 0.6 to 0.1; this tuning leads to improvements on the evaluation tasks, albeit requiring an increased number of inference steps.

The {\textsc BLT}-Entropy model demonstrates superior performance compared to the Llama 3 model in 4 of the 7 evaluated tasks, even though both were trained with an equivalent data volume measured in bytes. This enhanced performance can plausibly be attributed to two factors: (1) the more efficient allocation of training compute enabled by dynamic patching, and (2) the explicit modeling of information at the byte level rather than at the token level.

Conversely, {\textsc BLT}-Space yields poorer results than the Llama 3 tokenizer across all tasks except one, yet it notably decreases inference flops due to its increased average patch size of 6 bytes. In contrast, the bpe and entropy-patching models exhibit similar average patch sizes, both around 4.5 bytes, when evaluated on the training data mixture. Given an identical training budget, the model utilizing larger patches is able to process 30\% more data relative to the other two approaches, which could shift {\textsc BLT} farther from the point of compute-optimality.

\subsection{Patches Scale Better Than Tokens} In {\textsc BLT} architectures, it is possible to enlarge both the model and the patch size concurrently, without exceeding the original computational cost for training and inference or altering the fixed quantity of training data. The flexibility to freely adjust patch sizes distinguishes patch-based models, as this allows them to overcome the limitations in efficiency faced by traditional fixed-vocabulary token-based approaches, elaborated in Section~\ref{section:bpe-patch}. Employing larger patches results in reduced computational demands, permitting these savings to be redirected towards scaling up the global latent transformer, since it needs to be invoked less frequently.

A fixed inference scaling analysis is performed to evaluate whether increasing model size while reducing the number of steps on larger patches yields better performance compared to smaller models using more steps. We begin with 400m and 3.6B parameter models utilizing the Llama 2 tokenizer and identify flop-matched models employing the Llama 3 tokenizer and {\textsc BLT}-Entropy models configured with mean patch sizes of 6 and 8 bytes within the training datamix (refer to~\autoref{tab:fixed-inf-params} for specific model configurations). Models using an 8-byte patch size are assigned 3 encoder layers rather than 1. Each model is trained over different allocations of training flop budgets.

As illustrated in \autoref{fig:fixed_inference_scaling}, {\textsc BLT} models demonstrate more favorable scaling behavior than tokenization-based models across both categories of inference flops. While BPE models initially show stronger performance under limited training resources, their advantage diminishes beyond the compute-optimal threshold, at which point {\textsc BLT} models begin to excel. Consequently, allocating greater resources to pretraining can yield a model that outperforms alternatives when subject to the same inference constraints. For instance, within the class of 8B models—such as Llama 3.1—pretraining on data volumes approximately two orders of magnitude greater than the compute-optimal level for that scale exemplifies this approach.

The threshold at which {\textsc BLT} surpasses token-based models has moved slightly nearer to the compute-optimal point as we transition to models within higher flop classes (decreasing from 3x to 2.5x the compute-optimal expenditure). Additionally, for models with a patch size of 8, the scaling behavior in these larger flop categories is more pronounced, causing these models to outpace others earlier. As elaborated in Section~\ref{section:parameter-matched-scaling}, increasing patch size leads to performance more closely resembling that of BPE models at higher scales. This trend appears to result, in part, from the byte-level Encoder and Decoder accounting for a smaller fraction of total computational effort as scale increases, since their scaling is less rapid compared to that of the Latent Transformer. For instance, expanding the overall parameter count by a factor of 20, from 400M to 8B, results in only an approximate twofold increase in the local parameter count of {\textsc BLT}. Notably, alterations in patch size exclusively influence flops arising from the patch Latent Transformer, leaving the byte-level components unchanged. This factor explains why the {\textsc BLT}-Entropy ps=8 model ratio increased slightly, from 1.6x to 1.7x relative to Llama 2's size, when the model scale was increased.

To summarize, our investigation into patch-length scaling reveals that the {\textsc BLT} patch-based framework benefits from concurrent increases in both patch size and model scale, yielding superior scaling behavior. Notably, these favorable scaling properties appear to hold and potentially strengthen at even greater model capacities.

\section{Byte Modeling Improves Robustness} We further evaluate the robustness of {\textsc BLT} relative to token-based models that do not inherently encode byte-level detail, and we propose a strategy for integrating byte-level information into pretrained token-based architectures.
\label{sec:robustness}

\subsection{Character-Level Tasks}

An initial impetus for developing models operating at the byte level was their inherent resilience to noise occurring at this granular input scale, as well as their capacity to capture the internal structure of tokens—a capability with which existing models reliant on tokenizers often have difficulty. To quantitatively assess these attributes, we conduct supplementary evaluations using benchmarks specifically designed to test both the robustness of models to input-level noise and their sensitivity to character-level information across English and multiple languages, encompassing categories such as digits and phonemes. The outcomes of these assessments are summarized in Table~\ref{tab:char_tasks}.

\paragraph{Noisy Data} To assess and contrast the robustness of tokenizer-based models with {\textsc BLT}, we generate perturbed versions of the benchmark classification datasets detailed in Section~\ref{section:evals}. We implement five separate character-level perturbation techniques to introduce noise into the text: (a) \textit{AntSpeak}: This method renders the text in uppercase and separates each character with a space. (b) \textit{Drop}: This involves randomly eliminating 10\% of the text's characters. (c) \textit{RandomCase}: Here, half of the characters are randomly assigned to either uppercase or lowercase. (d) \textit{Repeat}: 20\% of the characters are selected and duplicated up to four times. (e) \textit{UpperCase}: This process transforms every character in the text to uppercase. For evaluation, each noising procedure is individually applied to the prompt, the completion, or both, creating distinct evaluation scenarios, and we present the mean performance. Results on the noised HellaSwag~\citep{Eval_hellaswag}, as summarized in Table \ref{tab:char_tasks}, indicate that {\textsc BLT} consistently exceeds the performance of tokenizer-based approaches in resilience to noise, achieving an average improvement of 8 points compared to a model trained under identical conditions, and surpasses the Llama 3.1 model, despite the latter being trained with significantly larger data.

\paragraph{Phonology - Grapheme-to-Phoneme (G2P)} We evaluate {\textsc BLT}'s performance in converting sequences of graphemes (written characters of a word) into their corresponding phonemic representations. Table \ref{tab:char_tasks} summarizes the outcomes for the G2P task under a 5-shot evaluation, utilizing the Phonology Bench~\citep{suvarna2024phonologybench}. The findings demonstrate that {\textsc BLT} surpasses the Llama 3 1T tokenizer-based baseline in this particular task.

\paragraph{CUTE}
To evaluate character-level capabilities, we benchmark {\textsc BLT} on the CUTE dataset~\citep{edman2024cute}, which consists of tasks that fall into three major types: compositional understanding, recognition of orthographic similarity, and sequence manipulation. This benchmark is notably difficult for typical tokenizer-based systems, which, although seemingly aware of token spellings, often fail to leverage this knowledge when asked to modify textual data. As reported in Table~\ref{tab:char_tasks}, {\textsc BLT}-Entropy achieves results exceeding both BPE Llama 3 variants by over 25 points. The model exhibits outstanding accuracy in character manipulation, scoring 99.9\% in both spelling-related tasks. These substantial gains, despite {\textsc BLT} being exposed to 16 times less training data than Llama 3.1, imply that BPE-based approaches struggle to internalize character-level patterns. Figure \ref{fig:cute_examples} provides example cases where the Llama 3 tokenizer model underperforms, yet our {\textsc BLT} model succeeds. The only instances where BPE surpasses are in word deletion and insertion tasks; while such manipulations present more difficulty for byte-level architectures, the performance gap is minor, and hierarchical composition from characters to words may ultimately prove more tractable than disassembling words into characters. For all CUTE tasks, we employ consistent evaluation procedures and retain the original Huggingface prompts. It should be noted that BPE models might obtain improved results with enhanced prompt engineering.

\paragraph{Low Resource Machine Translation} We assess the performance of {\textsc BLT} on both directions of translation across twenty-one low-resource languages, spanning six widely spoken language families and encompassing diverse scripts, utilizing the FLORES-101 benchmark~\citep{goyal2022flores}. The corresponding SentencePiece BLEU scores are presented in Table~\ref{tab:flores}. The findings indicate that {\textsc BLT} surpasses models using the Llama 3 tokenizer, offering a 2-point improvement when translating into English, and a 0.5-point gain when translating from English. In widely studied language pairs, the results of {\textsc BLT} are similar to or marginally higher than those of Llama 3. Nevertheless, for many language pairs within low-resource families, {\textsc BLT} distinctly outperforms Llama 3, highlighting the utility of byte-level modeling for capturing and generalizing rare or unique byte sequences found in the long tail of these languages.

\subsection{Training {\textsc BLT} from Llama 3}
\label{sec:distillation}
We investigate a procedure whereby {\textsc BLT} models can capitalize on pre-existing, pre-trained tokenizer-based architectures to facilitate more efficient and expedient training. Specifically, this is realized by initializing the global transformer parameters in {\textsc BLT} with those inherited from a pre-trained Llama 3.1 model. In the subsequent training phase, the global transformer parameters are fine-tuned using a learning rate set to one-tenth of that adopted for training the local encoder and decoder, maintaining the same number of optimization steps as used in Llama 3. A comparative analysis with a baseline {\textsc BLT} appears in Table~\ref{tab:distill}. The results indicate that {\textsc BLT} initialized from Llama 3.1 substantially outperforms both the original Llama 3 and the standard {\textsc BLT} models, all trained with an equivalent FLOPs budget. Furthermore, when evaluated against our {\textsc BLT}-Entropy model from Table~\ref{tab:evals}, which was exposed to a much larger corpus (1T tokens), the {\textsc BLT} model initialized from Llama 3.1 continues to demonstrate superior results on the MMLU benchmark, highlighting its efficacy in considerably lowering training FLOPs without sacrificing performance.

This approach can be interpreted as converting models that rely on tokenizers into ones that operate without tokenization, thereby adapting a pre-trained LLaMA 3.1 model into a {\textsc BLT} model. For thoroughness, we present results for the original LLaMA 3.1 model, which was trained on 15T tokens, in Table \ref{tab:distill}, and compare its performance to that of the LLaMA-based {\textsc BLT}. While our approach yields minor reductions in performance on MMLU and HumanEval, there are more pronounced declines observed on the remaining benchmarks. These outcomes indicate that additional research is necessary to better harness the capabilities of the pre-trained model, particularly by enhancing choices of data mixtures and tuning other hyperparameters.

\section{Ablations and Discussion}
\label{section:ablations}
Here, we present ablation studies that validate the architectural decisions behind {\textsc BLT}, as well as the selection of patching strategies and tuning of hyper-parameters for the {\textsc BLT} model with 8B parameters, trained on the {\textsc BLT}-1T corpus.

\paragraph{Entropy Model Hyper-parameters} In order to analyze how the size of the entropy model and the length of the context window influence scaling behavior, we train byte-level entropy transformer models spanning a range of sizes from 1m to 100m parameters, while also varying the context window length from 64 up to 512. For this investigation, we construct bpb versus training flop scaling curves using our $400m$ and $1b$ {\textsc BLT} models, all trained on the Llama-2 dataset, and display the results in \autoref{fig:entropy_ablation}. The analysis indicates that improvements in scaling are associated with increases in both the entropy model size and context window, although these gains begin to plateau as model sizes exceed 50m parameters.

\paragraph{Types of Patching}

We conduct an ablation study on the four patching methods presented in Section~\ref{section:patching}: (1) Strided Patching with strides set to 4 and 6, (2) Whitespace-based Patching, (3) BPE Tokenizer patching leveraging the Llama 3 tokenizer, and (4) Entropy-based patching utilizing a compact byte-level language model.

Although dynamic patching shortens the sequences, we ensure that the sequence length is regulated so that each patching approach has access to a comparable amount of context. Consequently, across all models, the expected number of bytes per sequence—during both training and inference—remains the same, minimizing confounds related to handling longer contexts. The findings from these ablation studies are illustrated in Figure~\ref{fig:scaling-compute-opt}. All patching strategies other than static patching demonstrate superior performance, with space patching closely matching the effectiveness of dynamic entropy-based patching.

In \autoref{tab:evals_ablation}, we report benchmark results for {\textsc BLT} models trained on the Llama 2 dataset, comparing standard tokenizer-based approaches, space patching, and entropy-based patching, all for a comparable number of training steps~\citep{dubey2024llama}. While space patching offers a more straightforward method that omits the need to evaluate an entropy model during training, our findings indicate that the improvements identified with entropy-based patching in scaling experiments (see Section~\ref{section:scaling}) also extend to performance on downstream benchmarks.\footnote{Space patching outcomes originate from earlier experiments without cross-attention; nevertheless, analogous patterns persist even when cross-attention is incorporated.}

\paragraph{Cross-Attention} 
In~\autoref{tab:crossattn_ablations_1b}, we systematically evaluate the effects of incorporating cross-attention at different stages within the encoder and decoder of {\textsc BLT}. For the encoder's cross-attention mechanism, we examine three strategies for initializing the queries: 1) assigning an identical learned embedding to all global states, 2) utilizing a hash embedding derived from the patch bytes, and 3) applying a pooling operation on the encoder's hidden representations of the patch bytes at the specified encoder layer.

Our results indicate that incorporating cross-attention within the \textit{decoder} yields the most significant benefits. While the encoder also experiences marginal gains from cross-attention, these improvements manifest primarily when queries are initialized via pooling. Moreover, the advantages of cross-attention are especially pronounced for the Common-Crawl dataset, with larger patch sizes enhancing these effects further.

\paragraph{n-gram Hash Embeddings} 

We experiment with n-gram hash embedding vocabulary sizes of 0, 100K, 200K, and 400K, with results summarized in Table~\ref{tab:ngram_ablations_1b}. Hash embeddings consistently improve performance across all evaluated domains, with especially notable improvements observed on Wikipedia and Github—showing a 0.04 bpb advantage, as opposed to a 0.01 bpb gain after 15k training steps at the 8B parameter scale. Furthermore, at the 8B scale, reducing the hash vocabulary from 500K to 300K yields a negligible change (~0.001 bpb) at 15k steps. These findings suggest that hash embeddings play a crucial role in enabling {\textsc BLT} to achieve parity with models that use conventional tokenizers; nevertheless, increases beyond 300K hash embeddings result in only marginal improvements. It is also evident that the benefits of hashing and cross-attention are largely additive, as each enhances model performance on different data distributions.

\paragraph{Local Model Hyperparameters} Table \ref{tab:local_layers} presents an ablation study exploring different configurations for the layer count in both the local encoder and decoder. Our findings indicate that, particularly when hash n-gram embeddings are utilized, {\textsc BLT} achieves strong performance with a very shallow encoder—consisting of only a single layer—while the decoder benefits from a greater depth.

\section{Related Work}
\label{section:related_work}

\textbf{Character-Level RNNs:} The task of Character Language Modeling has maintained its prominence since the initial development of neural architectures~\citep{sutskever2011generating,mikolov2012subword,graves2013generating}, primarily due to their inherent ability to seamlessly represent out-of-vocabulary terms, eliminating the reliance on back-off strategies. In particular, \cite{kim2016character} develop a model that exclusively utilizes characters on the input side, leveraging convolutional and highway networks, which subsequently feed into LSTM-based RNNs. This approach not only achieves parity with the then state-of-the-art RNN language models in English, but also surpasses them on languages with greater morphological complexity—underscoring a key benefit of character-level LLMs. Additionally, \cite{kenter2018byte} employ byte-level LSTM architectures for machine comprehension, demonstrating superior performance compared to word-level counterparts, notably in Turkish and Russian, both highly inflectional languages. Similarly, \cite{zhang2015character} adopt character-based convolutional models for classification, observing improved results over word-level models in certain applications. Building further, \citet{chung2019hierarchical} implement hierarchical LSTMs with boundary-detection mechanisms at every layer, enabling the discovery of latent textual hierarchies and advancing character-level language modeling effectiveness. ByteNet, proposed by \cite{kalchbrenner2016neural}, introduces convolutional network layers applied to character sequences, in contrast to attention-based mechanisms, specifically for machine translation tasks.

\textbf{Character-Level Transformers:} The introduction of attention mechanisms in transformer architectures~\citep{vaswani2017attention}, along with the adoption of subword tokenization techniques~\citep{sennrich-etal-2016-neural}, resulted in substantial advancements in neural models for language modeling and related benchmarks. Nevertheless, the reliance on word and sub-word tokens inherently establishes a preferred abstraction layer for these models. Seeking to bridge transformer architectures with promising early outcomes in character-level language modeling, \citet{al2019character} implemented very deep transformers augmented with auxiliary loss functions, training transformer-based models that achieved superior results relative to earlier LSTM-based character language models. Despite these improvements, the models continued to lag noticeably behind their word-level LLM counterparts. Similarly, GPT-2~\citep{radfordgpt2} found that on large-scale datasets such as the 1 Billion Word Benchmark, byte-level language models were still not able to match the performance of word-level models.

\cite{Choe2019BridgingTG} showed that transformer-based large language models (LLMs) operating at the byte level can surpass subword-level LLMs with similar parameter counts; however, these models demand significantly more computational resources and entail longer training times. Along similar lines, \cite{el2020characterbert} introduced a BERT variant, CharFormer, which constructs word-level representations through convolutions applied to character embeddings. While this approach yields better performance within the medical domain, it also incurs higher computational costs. The CANINE model, developed by \cite{clark2022canine}, features 150M parameters in an encoder-only architecture that processes character sequences directly. CANINE’s backbone consists of a deep transformer stack, akin to the global module in our work, complemented by a local transformer and strided convolutions to reduce input dimensionality, thereby outperforming a comparable token-based encoder (mBERT) on multilingual downstream tasks. ByT5~\citep{xue2022byt5} proposed byte-level encoder-decoder models that forego patching strategies; while they achieve superior noise robustness and rival token-based models with four times less training data, omitting patching necessitates computationally expensive attention across all input bytes, leading to high compute demands. Modeling data directly at the byte level increases input sequence lengths, introducing efficiency bottlenecks when scaling such architectures. More recently, leveraging the Mamba Architecture~\citep{gu2023mamba}—which supports a fixed-size memory state over extensive context windows—\cite{wang2024mambabyte} presented a byte-level Mamba-based model also without patching. At a 350M parameter scale, their model surpasses byte-level transformer approaches in flop-matched scenarios in terms of bits-per-byte across various datasets.

\textbf{Patching-based approaches:} Leveraging patching strategies can help substantially reduce the computational overhead—specifically, the number of flops—required by byte-level LLMs, all while maintaining comparable model performance. A number of studies have provided promising results on limited scales, both in model capacity and data size. \cite{nawrot-etal-2022-hierarchical} introduce a static patching methodology involving downsampling and upsampling mechanisms and propose the hourglass transformer architecture, which surpasses competing byte-level models at the 150M parameter scale. Building upon this, \cite{nawrot-etal-2023-efficient} propose several dynamic patching methods: these include an end-to-end trainable boundary predictor, a variant guided by specific tokenizers, and an entropy-based technique akin to {\textsc BLT}. Their results indicate that these dynamic patching techniques outperform standard transformers on language modeling benchmarks when applied to models with around 40M parameters trained on roughly 400M tokens. Separately, \cite{lester2024training} explore the use of arithmetic coding to compress input sequences, enabling greater compression than standard BPE. Through their equal-info windows approach, they demonstrate improvements over byte-level baselines in language modeling tasks, though results fall short of those achieved by subword-based models.

This study is heavily informed by MegaByte~\citep{yu2023megabyte}, a causal decoder-only LLM that transforms bytes into patches using a predetermined fixed patching strategy and merges their representations through concatenation. On the decoding side, MegaByte applies a localized model to reconstruct the original bytes from patches. MegaByte achieves parity with tokenizer-based models possessing 1B parameters on datasets roughly 400B bytes in size. In all of our analyses, we compare directly with MegaByte, revealing that static patching methods underperform relative to state-of-the-art tokenizer-based architectures, particularly under fixed compute budgets; we further show how {\textsc BLT} narrows this performance gap. Likewise, \cite{slagle2024spacebyte} note similar shortcomings in MegaByte, proposing enhancements such as dynamic patching—using whitespace and similar characters as segmentation points—and incorporating a local encoder. They report that these modifications yield performance gains over tokenizer-based transformers on certain datasets, including Github and arXiv, also at the 1B parameter level. Our own experimentation with this approach corroborates their findings but also indicates that additional design advancements are essential to bring byte-level models to the performance levels of the most recent state-of-the-art token-based systems like Llama 3.

\section{Limitations and Future Work}

In this study, our architectural decisions are informed by training models using the optimal number of steps identified for Llama~3~\citep{dubey2024llama}. Nevertheless, these scaling laws were derived based on BPE-level transformers, which might not yield ideal data-to-parameter ratios for the {\textsc BLT} model. Determining scaling laws specifically suited to {\textsc BLT}—which could result in more advantageous scaling properties—remains a subject for future investigation. Furthermore, as the bulk of our experiments were limited to model sizes up to 1B parameters, it remains plausible that optimal architectural configurations may differ when extended to models with 8B parameters or greater, potentially leading to enhanced performance at larger model sizes.

Current transformer libraries and code repositories are primarily optimized for high efficiency with tokenizer-based transformer models. Although our experiments are theoretically matched in terms of FLOPs and we employ optimized solutions like FlexAttention for non-standard transformer layers, the actual wall-clock performance of our implementations may lag behind those of tokenizer-based counterparts. Consequently, there remains potential for additional optimization to bridge this gap.

Whereas {\textsc BLT} relies on a separately trained entropy model for its patching process, jointly optimizing the patching model in an end-to-end manner offers a promising avenue for future research. In Section \ref{sec:distillation}, we provide preliminary experimental results that demonstrate early success in ``byte-ifying'' tokenizer-based models like Llama 3—which have been trained on datasets exceeding 10T tokens—by initializing and freezing the global transformer with pre-existing weights. Further exploration along this trajectory could yield techniques that preserve the advantages of bytefying while potentially surpassing the original tokenizer-based models' performance, all without the need for complete retraining.

\section{Conclusion}
\label{section:conclusion}

In this work, we introduce the Byte Latent Transformer (\textbf{BLT}), an architecture that challenges the standard reliance on fixed-vocabulary tokenization in large language models. By employing a trainable, adaptive strategy for aggregating bytes into variable-length patches, {\textsc BLT} dynamically adjusts computational allocation in response to the underlying data complexity. This adaptive processing yields notable gains in computational efficiency and robustness. Through comprehensive scaling experiments, we show that {\textsc BLT} models can achieve parity with tokenization-based approaches such as Llama 3, across scales reaching 8B parameters and 4T bytes of data, while offering up to 50\% savings in inference computational costs with only marginal decreases in evaluation performance. Additionally, {\textsc BLT} facilitates a novel scaling direction, enabling concurrent scaling of both model and patch sizes without increasing overall inference requirements—a capability particularly well-suited for real-world compute constraints. By operating directly on raw byte representations, {\textsc BLT} enhances a model’s capacity to manage rare or outlier data, demonstrating improved robustness to noisy inputs and finer-grained comprehension of sub-word elements. Collectively, these findings highlight {\textsc BLT} as a compelling alternative to classical tokenization-driven models, delivering an adaptable and efficient foundation for next-generation language modeling.

\end{document}